\appendix

\section{Unités, invariance, et ce qui est calculable}
\label{app:units}

\paragraph{Le nat.} L'entropie $H=-\sum_i p_i\log p_i$ n'est définie qu'une fois
la base du logarithme choisie, et ce choix \emph{est} le choix de l'unité. Avec
$\log_2$ on obtient des bits (shannons), avec $\log_{10}$ des dits, et avec le
logarithme népérien des \emph{nats}. Un nat est l'information apportée par une
observation qui divise l'espace des possibles par $e\approx2{,}718$, là où un bit
le divise par deux. La conversion est
$1~\text{nat}=1/\log 2\approx1{,}4427$~bit. Tout ce papier emploie le logarithme
népérien : entropies, déficits, divergences, rayons de scénario et indice
structurel sont en nats.

Trois lectures opérationnelles, de la plus utile à la plus abstraite. D'abord le
taux de change établi en section~\ref{sec:stress} : un scénario purement
directionnel de $\eta$ nats est un mouvement de $\sqrt{2\eta}$ sigmas, donc
$k$~sigmas coûtent $k^2/2$ nats. Ensuite le codage : $\KL{p}{q}$ nats est
l'excès de longueur de code attendu lorsqu'on encode des données de loi $p$ avec
un code optimal pour $q$. Enfin le test d'hypothèse : une divergence de un nat
par observation signifie que le rapport de vraisemblance en faveur de $p$ croît
en moyenne d'un facteur $e$ à chaque observation.

\paragraph{$h$ dépend des unités, les divergences non.} L'entropie différentielle
d'une variable continue vérifie, pour toute transformation affine inversible,
\begin{equation}
\label{eq:h-scaling}
h(aX+b) \;=\; h(X) + \log|a| ,
\qquad\text{et en dimension } N,\quad h(AX+b)=h(X)+\log|\det A| .
\end{equation}
Exprimer les rendements en pourcentages plutôt qu'en décimales ajoute donc
$N\log 100$ à $h$ : le \emph{niveau} de $h$, et par conséquent celui de
$H^{vol}$ et de $H^{cov}$, n'a aucune valeur absolue interprétable. En revanche
$\KL{p}{q}$, $I(X;Y)$, $\DD$ et $\JJ$ sont invariants par tout changement de
variable inversible, les jacobiens se simplifiant entre numérateur et
dénominateur du logarithme. \emph{Un chiffre en nats n'est comparable d'un jeu
de données à l'autre que si c'est une divergence.} C'est la raison structurelle
pour laquelle l'article construit un indice sans échelle.

\paragraph{Comment calcule-t-on $h(X)$ ?} Dans le cas gaussien, en forme fermée
(annexe~\ref{app:gauss}). Dans le cas général, on ne le peut pas sans connaître
$p$. Des estimateurs non paramétriques existent --- espacements de
\citet{vasicek1976}, plus proches voisins de Kozachenko--Leonenko, méthodes à
noyau --- mais tous se dégradent rapidement au-delà de trois ou quatre
dimensions. Estimer une entropie différentielle en dimension~20 sur une fenêtre
de 252~jours est hors de portée.

L'article ne calcule donc \emph{jamais} $h(p)$. Il calcule $h(q)$ en forme
fermée à partir de $\widehat\Sigma$, et le déficit $\DD$ par une approximation
de la \emph{différence} d'entropies marginales. Une différence d'entropies est
sans dimension et estimable ; une entropie différentielle nue ne l'est pas.
Toute l'architecture --- décomposer plutôt qu'estimer directement --- découle de
cette contrainte.

\section{Entropie gaussienne : démonstration de l'équation~\eqref{eq:ceiling}}
\label{app:gauss}

\begin{proof}[Étape A : entropie d'une gaussienne multivariée]
Pour $q=\mathcal N(\mu,\Sigma)$ sur $\mathbb{R}^N$,
\[
\log q(x) = -\tfrac12\left[N\log 2\pi + \log\det\Sigma + (x-\mu)'\Sigma^{-1}(x-\mu)\right],
\]
donc
\[
h(q) = -\mathbb{E}_q[\log q] = \tfrac12 N\log 2\pi + \tfrac12\log\det\Sigma
+ \tfrac12\,\mathbb{E}\!\left[(X-\mu)'\Sigma^{-1}(X-\mu)\right].
\]
La dernière espérance se calcule par l'astuce de la trace : un scalaire égale sa
trace, et la trace est cyclique, d'où
\[
\mathbb{E}\!\left[(X-\mu)'\Sigma^{-1}(X-\mu)\right]
= \mathbb{E}\!\left[\operatorname{tr}\!\left(\Sigma^{-1}(X-\mu)(X-\mu)'\right)\right]
= \operatorname{tr}(\Sigma^{-1}\Sigma) = \operatorname{tr}(I_N) = N .
\]
En regroupant, $h(q)=\tfrac12\left[N\log2\pi+\log\det\Sigma+N\right]
=\tfrac12\log\!\left((2\pi e)^N\det\Sigma\right)$.
\end{proof}

\begin{proof}[Étape B : factorisation du déterminant]
Avec $\Sigma=DRD$, $D=\operatorname{diag}(\sigma_1,\dots,\sigma_N)$ et $R$ la
matrice de corrélation, la multiplicativité du déterminant donne
$\det\Sigma=\det(D)\det(R)\det(D)=\left(\prod_i\sigma_i\right)^2\det R$.
\end{proof}

\begin{proof}[Étape C : substitution]
$\tfrac12\log\det\Sigma
=\tfrac12\log\!\left[\left(\prod_i\sigma_i\right)^2\det R\right]
=\sum_i\log\sigma_i+\tfrac12\log\det R$, d'où
$h(q)=\tfrac N2\log(2\pi e)+\sum_i\log\sigma_i+\tfrac12\log\det R$, qui est
l'équation~\eqref{eq:ceiling}.
\end{proof}

\paragraph{Signe du canal de dépendance.} L'inégalité de Hadamard donne
$\det R\le\prod_i R_{ii}=1$ pour une matrice semi-définie positive. Une
démonstration plus parlante passe par le spectre : $\det R=\prod_k\lambda_k$
avec $\sum_k\lambda_k=\operatorname{tr}R=N$, et l'inégalité arithmético-géométrique
donne $\prod_k\lambda_k\le\left(\sum_k\lambda_k/N\right)^N=1$, avec égalité si
et seulement si toutes les valeurs propres valent~1, c'est-à-dire $R=I$. Donc
$H^{dep}\le0$ : \emph{la corrélation détruit toujours de l'entropie}.

Dans le cas bivarié, $\det R=1-\rho^2$ et $H^{dep}=\tfrac12\log(1-\rho^2)$ ; pour
$\rho=0{,}9$ cela vaut $-0{,}83$~nat. Remarquons que $-H^{dep}$ est alors
exactement l'information mutuelle entre les deux actifs : la corrélation est de
l'information partagée, donc de l'entropie détruite.

\paragraph{Note de lecture sur les tableaux.} Les colonnes $H^{cov}$ des
tableaux de résultats de la version arXiv~v1 rapportent $H^{vol}+H^{dep}$ et
\textbf{omettent la constante} $\tfrac N2\log(2\pi e)$, qui vaut $28{,}38$~nats
pour $N=20$. Ce décalage constant est sans conséquence puisque seules les
différences entre régimes et les variations temporelles sont interprétées, mais
il explique que les niveaux rapportés soient négatifs.

\section{Le déficit d'entropie : pourquoi $\DD=h(q)-h(p)$}
\label{app:deficit}

La divergence de Kullback--Leibler entre une loi $p$ et une loi de référence $q$
est $\KL{p}{q}=\int p\log(p/q)$, positive et nulle si et seulement si $p=q$
presque partout. Elle n'est ni symétrique ni munie d'une inégalité triangulaire :
ce n'est pas une distance, d'où le terme de divergence. Les lettres \og KL \fg{}
sont les initiales de Kullback et Leibler.

\begin{proposition}
\label{prop:deficit}
Si $q=\mathcal N(\mu,\Sigma)$ a \emph{les deux premiers moments de $p$}, alors
$\KL{p}{q}=h(q)-h(p)$.
\end{proposition}

\begin{proof}
$\KL{p}{q}=-h(p)-\mathbb{E}_p[\log q(X)]$. Or $\log q$ est exactement
quadratique en $x$, donc son espérance ne dépend que de la moyenne et de la
covariance de la loi sous laquelle on l'évalue. Comme $p$ et $q$ partagent ces
deux moments,
$\mathbb{E}_p\!\left[(X-\mu)'\Sigma^{-1}(X-\mu)\right]=\operatorname{tr}(\Sigma^{-1}\Sigma)=N
=\mathbb{E}_q[\cdot]$, d'où $\mathbb{E}_p[\log q]=\mathbb{E}_q[\log q]=-h(q)$ et
donc $\KL{p}{q}=h(q)-h(p)$.
\end{proof}

L'hypothèse de coïncidence des moments est essentielle : sans elle l'identité
tombe. C'est tout le contenu de la proposition~\ref{prop:maxent}, lue comme une
mesure plutôt que comme un théorème d'existence. Économiquement, $\DD$ est le
prix en nats du mensonge \og le marché est gaussien \fg{} : l'incertitude
apparente que la description gaussienne surestime, parce que la loi véritable est
plus concentrée que ne le laisse croire sa seule covariance.

\section{L'indice structurel est une divergence}
\label{app:index}

\begin{proposition}
\label{prop:index}
Avec $q_i=\mathcal N(\mu_i,\sigma_i^2)$ les marginales gaussiennes de mêmes
échelles que $p$,
\[
\KL{p}{\textstyle\prod_i q_i} \;=\; \DD - H^{dep} \;=\; \JJ .
\]
\end{proposition}

\begin{proof}
Par définition,
\[
\begin{aligned}
\KL{p}{\textstyle\prod_i q_i} &= -h(p) - \mathbb{E}_p\!\left[\textstyle\sum_i\log q_i(X_i)\right]
= -h(p) + \sum_i\left[\tfrac12\log(2\pi\sigma_i^2)+\tfrac12\right]\\
&= -h(p) + \sum_i\log\sigma_i + \tfrac N2\log(2\pi e),
\end{aligned}
\]
où l'on a utilisé $\mathbb{E}_p[(X_i-\mu_i)^2]=\sigma_i^2$, vrai par
construction. En y injectant l'équation~\eqref{eq:three}, soit
$h(p)=\tfrac N2\log(2\pi e)+\sum_i\log\sigma_i+H^{dep}-\DD$, les termes
d'échelle et la constante se télescopent et il reste $\DD-H^{dep}$.
\end{proof}

C'est une identité exacte, non une approximation, et c'est elle qui justifie la
définition~\ref{def:J}. Retirer le canal d'échelle de \eqref{eq:three} pourrait
passer pour un bricolage ; l'objet qui reste se trouve être une divergence de
Kullback--Leibler en bonne et due forme, et hérite donc gratuitement de trois
propriétés : elle est positive ; elle est nulle si et seulement si les rendements
sont indépendants \emph{et} gaussiens, ce qui fournit un point de référence bien
défini plutôt qu'arbitraire ; et elle est invariante par changement d'unités
actif par actif, puisque $p$ et $\prod_i q_i$ se transforment de la même façon.

Il est utile de rapprocher cette identité de la règle de chaînage générale
$\KL{p}{\prod_i q_i}=\sum_i\KL{p_i}{q_i}+I(x)$, où
$I(x)=\KL{p}{\prod_i p_i}$ est la multi-information. Le premier terme est la
somme des déficits marginaux $\sum_i\Dun(x_i)$, le second la dépendance totale.
L'estimateur employé dans l'article estime le premier terme par l'approximation
de \citet{hyvarinen1998} et remplace le second par sa part linéaire
$-H^{dep}$ ; la dépendance de queue non linéaire est donc omise de $\JJ$ et
mesurée séparément par un taux de co-dépassement.

\section{Information mutuelle : les deux propositions et un exemple}
\label{app:mi}

\paragraph{Démonstration de la proposition~\ref{prop:conditioning}.} Par
définition de l'entropie conditionnelle,
$\mathbb{E}_Y[h(X\mid Y)]=h(X,Y)-h(Y)$, et par
\eqref{eq:mi}, $I(X;Y)=h(X)+h(Y)-h(X,Y)$. En sommant les deux,
$\mathbb{E}_Y[h(X\mid Y)]=h(X)-I(X;Y)$. La positivité de $I$ découle de celle de
la divergence de Kullback--Leibler dont elle est un cas particulier, avec
égalité si et seulement si $p_{XY}=p_X\otimes p_Y$.

\paragraph{Démonstration de la proposition~\ref{prop:constraint}.} Immédiate :
une contrainte supplémentaire ne peut que réduire l'ensemble réalisable, et le
supremum d'une fonctionnelle sur un ensemble plus petit ne peut pas augmenter.

\paragraph{Pourquoi les deux ne sont pas équivalents.} La première quantifie
mais ne vaut qu'en moyenne ; la seconde vaut réalisation par réalisation mais
porte sur le plafond et ne quantifie rien. Un contre-exemple discret suffit à
montrer qu'une réalisation particulière peut \emph{augmenter} l'entropie. Soit
$X$ binaire avec $\mathbb{P}(X=a)=0{,}99$, d'entropie $H(X)=0{,}056$~nat. Soit
un signal $Y$ qui, avec probabilité $0{,}98$, lève complètement le doute
($H(X\mid Y=0)=0$), et avec probabilité $0{,}02$ conduit à
$\mathbb{P}(X=a\mid Y=1)=1/2$, soit $H(X\mid Y=1)=\log 2=0{,}693$~nat --- douze
fois l'entropie a priori. La moyenne vaut néanmoins
$0{,}98\times0+0{,}02\times0{,}693=0{,}014$~nat, bien en dessous de $0{,}056$,
et $I(X;Y)=0{,}042$~nat. C'est exactement la raison pour laquelle l'article
prédit une composante de saut \emph{principalement}, et non exclusivement,
unilatérale.

\paragraph{Exemple gaussien chiffré.} Soit $X\sim\mathcal N(0,\sigma^2)$ le
rendement de demain et $Y=X+\varepsilon$ un signal bruité,
$\varepsilon\sim\mathcal N(0,\sigma_\varepsilon^2)$ indépendant. Alors
\begin{equation}
\label{eq:mi-gauss}
I(X;Y)=\tfrac12\log\!\left(1+\frac{\sigma^2}{\sigma_\varepsilon^2}\right),
\qquad
\operatorname{Var}(X\mid Y)=\frac{\sigma^2\sigma_\varepsilon^2}{\sigma^2+\sigma_\varepsilon^2},
\end{equation}
la variance a posteriori ne dépendant pas de la réalisation $y$. Avec
$\sigma=1\%$ par jour :

\begin{center}
\small
\begin{tabular}{@{}lrrrr@{}}
\toprule
Qualité du signal & $\sigma_\varepsilon$ & $I(X;Y)$ (nats) & vol.\ a posteriori & équivalent $\sqrt{2\eta}$ \\
\midrule
tranchant   & $0{,}5\%$ & $0{,}805$ & $0{,}45\%$ & $1{,}27\sigma$ \\
moyen       & $1{,}0\%$ & $0{,}347$ & $0{,}71\%$ & $0{,}83\sigma$ \\
vague       & $3{,}0\%$ & $0{,}053$ & $0{,}95\%$ & $0{,}32\sigma$ \\
\bottomrule
\end{tabular}
\end{center}

Un facteur quinze sépare la taille du saut d'entropie provoqué par une nouvelle
qui tranche de celui provoqué par une rumeur. La dernière colonne traduit chaque
saut dans l'échelle de sévérité de la section~\ref{sec:stress} : c'est le
mouvement directionnel qu'un scénario de même rayon produirait.

